Let $f:\mathbb N\to\mathbb R$ be additive on coprime factors and suppose $|f(n+1)-f(n)|<C$ for all $n$. Must $f(n)=c\log n+O(1)$ for some real $c$? Wirsing proved that the answer is yes. This document proves an exact characterization of the possible bounded additive remainders, verifies that they give the required bounded increments, and shows why the remainder cannot be omitted. It does not derive the existence of $c$ from the increment hypothesis alone.
